 \documentclass[letterpaper]{article} % DO NOT CHANGE THIS
\usepackage[submission]{aaai2027}  % DO NOT CHANGE THIS
% The serif, sans-serif, and monospaced fonts are loaded automatically by
% aaai2027.sty (newtxtext, helvet, courier). DO NOT add \usepackage{times},
% \usepackage{helvet}, \usepackage{courier}, or any other font package.
\usepackage[hyphens]{url}  % DO NOT CHANGE THIS
\usepackage{graphicx} % DO NOT CHANGE THIS
\urlstyle{rm} % DO NOT CHANGE THIS
\def\UrlFont{\rm}  % DO NOT CHANGE THIS
\usepackage{natbib}  % DO NOT CHANGE THIS AND DO NOT ADD ANY OPTIONS TO IT
\usepackage{caption} % DO NOT CHANGE THIS AND DO NOT ADD ANY OPTIONS TO IT
\frenchspacing  % DO NOT CHANGE THIS
%
% These are recommended to typeset algorithms but not required. See the subsubsection on algorithms. Remove them if you don't have algorithms in your paper.
\usepackage{algorithm}
\usepackage{algorithmic}

%
% These are recommended to typeset listings but not required. See the subsubsection on listing. Remove this block if you don't have listings in your paper.
\usepackage{newfloat}
\usepackage{listings}
\DeclareCaptionStyle{ruled}{labelfont=normalfont,labelsep=colon,strut=off} % DO NOT CHANGE THIS
\lstset{%
	basicstyle={\footnotesize\ttfamily},% footnotesize acceptable for monospace
	numbers=left,numberstyle=\footnotesize,xleftmargin=2em,% show line numbers, remove this entire line if you don't want the numbers.
	aboveskip=0pt,belowskip=0pt,%
	showstringspaces=false,tabsize=2,breaklines=true}
\floatstyle{ruled}
\newfloat{listing}{tb}{lst}{}
\floatname{listing}{Listing}

%
% Recommended for better-looking tables
\usepackage{booktabs}

% Math and theorem environments (permitted alongside the AAAI kit).
\usepackage{amsmath}
\usepackage{amssymb}
\usepackage{amsthm}
\theoremstyle{definition} % upright body: our propositions are prose-heavy
\newtheorem{proposition}{Proposition}
\newtheorem*{remark}{Remark} % 한글판 "비고"와 같이 무번호

% Notation macros (keep in sync with paper-ko.md).
\newcommand{\lval}{\ell_{\mathrm{val}}} % mean loss on the server-held validation set
\usepackage{kotex}
\usepackage[dvipsnames]{xcolor}
\newcommand{\yhha}[1]
{{\leavevmode\color{SkyBlue}[yhha:#1]}}
\newcommand{\jwjung}[1]
{{\leavevmode\color{Green}[jwjung:#1]}}

%
% Keep the \pdfinfo as shown here. There's no need
% for you to add the /Title and /Author tags.
\pdfinfo{
/TemplateVersion (2027.1)
}

% DISALLOWED PACKAGES
% \usepackage{authblk} -- This package is specifically forbidden
% \usepackage{balance} -- This package is specifically forbidden
% \usepackage{CJK} -- This package is specifically forbidden
% \usepackage{float} -- This package is specifically forbidden
% \usepackage{flushend} -- This package is specifically forbidden
% \usepackage{fullpage} -- This package is specifically forbidden
% \usepackage{geometry} -- This package is specifically forbidden
% \usepackage{hyperref} -- This package is specifically forbidden
% \usepackage{navigator} -- This package is specifically forbidden
% (or any other package that embeds links such as navigator or hyperref)
% \usepackage{indentfirst} -- This package is specifically forbidden
% \usepackage{layout} -- This package is specifically forbidden
% \usepackage{multicol} -- This package is specifically forbidden
% \usepackage{nameref} -- This package is specifically forbidden
% \usepackage{savetrees} -- This package is specifically forbidden
% \usepackage{setspace} -- This package is specifically forbidden
% \usepackage{stfloats} -- This package is specifically forbidden
% \usepackage{tabu} -- This package is specifically forbidden
% \usepackage{titlesec} -- This package is specifically forbidden
% \usepackage{tocbibind} -- This package is specifically forbidden
% \usepackage{ulem} -- This package is specifically forbidden
% \usepackage{wrapfig} -- This package is specifically forbidden
% DISALLOWED COMMANDS
% \nocopyright -- Your paper will not be published if you use this command
% \addtolength -- This command may not be used
% \balance -- This command may not be used
% \baselinestretch -- Your paper will not be published if you use this command
% \clearpage -- No page breaks of any kind may be used for the final version of your paper
% \columnsep -- This command may not be used
% \newpage -- No page breaks of any kind may be used for the final version of your paper
% \pagebreak -- No page breaks of any kind may be used for the final version of your paper
% \pagestyle -- This command may not be used
% \tiny -- This is not an acceptable font size.
% \vspace{- -- No negative value may be used in proximity of a caption, figure, table, section, subsection, subsubsection, or reference
% \vskip{- -- No negative value may be used to alter spacing above or below a caption, figure, table, section, subsection, subsubsection, or reference

\setcounter{secnumdepth}{2} % 0에서 2로 변경(kit 허용): paper-ko.md의 §x.y 상호참조 스타일에 맞춰 섹션 번호 표시

% The file aaai2027.sty is the style file for AAAI Press
% proceedings, working notes, and technical reports.
%

% Title
% Your title must be in mixed case, not sentence case (AAAI 규정).
% [작업 제목] paper-ko.md 제목의 영문판. "Exact-Oracle Fidelity"는 dual-oracle 표현 정리(retrain
% 특성화 강등, §6.2) 이후 개명 검토가 보류 중 — 확정되면 여기와 paper-ko.md를 함께 갱신할 것.
\title{Flirds: Efficient Client-Level Contribution Evaluation in Federated Learning via In-Run Data Shapley}
\author{ 
    % 익명 제출 지침: 저자는 "Anonymous Submission" 단독, 소속은 비워 둔다.
    Anonymous Submission
}
\affiliations{
}

\begin{document}

\maketitle

\begin{abstract}
Client contribution evaluation in federated learning (FL) quantifies how much
each participant's update contributed to improving the jointly trained model.
Because an FL server does not hold the clients' raw data, it cannot by itself
perform the subset retraining on which classical data valuation is based, and
contributions must instead be evaluated from the updates received in actual
rounds. The Shapley value is widely used for contribution evaluation because
it provides a consistent attribution principle within a specified cooperative
game, but computing it exactly requires a number of coalition evaluations
exponential in the number of participating clients. Existing approximations
reduce the number of evaluations through sampling and truncation but still
require repeated evaluation of coalition models, and some cheaper alternative
scores and influence scores do not directly approximate the Shapley value of
a stated game, or compute values defined by a different utility or weighting.
What is needed, therefore, is a client-level contribution evaluation method
that is explicit not only about computational efficiency but also about what
it approximates and how its error should be interpreted.
To this end we propose \textbf{Flirds} (Federated Learning In-Run Data
Shapley), a client-level contribution estimator. Flirds defines a cooperative
game whose players are the client updates actually received by the server,
with the round's server aggregation weights held fixed across all coalitions;
the update of the grand coalition therefore coincides exactly with the
aggregation the server actually performed. We take as our reference the
per-round Shapley values obtained by exhaustively evaluating this game over
all coalitions, and extend the Taylor calculation of In-Run Data Shapley (IRDS)
to client updates, computing the Shapley value of a second-order surrogate
game in closed form. This replaces the exponential coalition evaluations with
one server-side Hessian-vector product per round and one inner product per
participant, with no additional communication or client-side computation.
Moreover, because the reference and the estimate target the same game, their
difference can be interpreted as the Taylor residual rather than a difference
in game definition or sampling error. Experiments validate the accuracy of
the estimates through ranking and value comparisons against the exhaustively
enumerated Shapley values, and their practical effectiveness through actual
training improvement under a shared intervention policy together with
measured computational cost.
\end{abstract}

% Uncomment the following to link to your code, datasets, an extended version or similar.
% You must keep this block between (not within) the abstract and the main body of the paper.
% Make sure that you do not de-anonymize yourself with these links.
% \begin{links}
%     \link{Code}{https://aaai.org/example/code}
%     \link{Datasets}{https://aaai.org/example/datasets}
%     \link{Extended version}{https://aaai.org/example/extended-version}
% \end{links}

\section{Introduction}
\label{sec:intro}

In federated learning (FL), clients jointly train a single model without
sharing their raw data \citep{mcmahan2016federated}. In each round, the
server collects the updates sent by the participating clients and improves
the model, but the aggregate alone does not directly reveal how much each
client contributed to that improvement. Such contribution estimates can serve
as grounds not only for participant compensation but also for data-quality
management, free-rider detection, client selection, and debugging of the
training process. The Shapley value from cooperative game theory provides a
principle for consistently allocating a specified utility among the
participants, and has therefore been widely used for evaluating data and
client contributions \citep{shapley1953value,ghorbani2019data}.

Classical Data Shapley retrains the model from scratch with each subset of
clients and compares the resulting final performance. An FL server, however,
does not hold the clients' raw data, so it cannot perform subset retraining
on its own; retraining amounts to re-running the entire federation and
repeatedly requesting the clients' cooperation. This definition is therefore
not an option available to the server for reasons of information structure,
before cost even enters, and still less can it serve as an online evaluation
tool in settings where clients participate intermittently. Prior federated
Shapley methods have instead taken the updates the server receives in an
actual round as the units of evaluation, computed per-round contributions,
and accumulated them over training. Even then, obtaining exact Shapley values
requires enumerating participant subsets in every round, reconstructing a
model for each coalition, and evaluating it on a validation set, so the cost
grows exponentially with the number of participants. Sampling, truncation,
interpolation, or separate alternative scores can reduce this cost, but they
may add sampling variability or compute a utility different from the original
target, which makes estimation error hard to distinguish from a change in the
game definition.

The problem this paper addresses is therefore the following: using only the
individual updates observed by the server and a validation set, can we
attribute the validation-loss improvement produced by an actual FedAvg round
to the participating clients in the Shapley sense, while avoiding
per-coalition model evaluation? This requires first defining a game whose
grand coalition coincides with the aggregation the server actually performed,
and that game must be computable online under partial participation. It must
also be possible to state separately which game's Shapley value the computed
quantity approximates and where the approximation error arises.

We propose \textbf{Flirds} (Federated Learning In-Run Data Shapley), a
client-level contribution estimator for this problem. Flirds views each
client's weighted update in an actual round as a player and defines the round
utility as the decrease in validation loss obtained when only a subset of
those updates is applied. The weight of each client is not renormalized
within each coalition but held fixed at the value used in the actual server
aggregation. The model displacement of the grand coalition therefore
coincides exactly with the actual FedAvg update. Exhaustively evaluating this
round game yields the per-round Shapley values, but still requires
exponentially many validation evaluations per round. Flirds removes this
enumeration by using the computational idea of In-Run Data Shapley (IRDS): it
computes in closed form the Shapley value of the second-order Taylor
surrogate of the validation loss.

This construction connects the method directly to the federated learning
setting. The principal additional server computation Flirds requires per
round is one Hessian-vector product (HVP) on the validation set, and the
per-participant computation is a single inner product. It needs no additional
communication or client-side computation, and the values of partially
participating rounds can be accumulated online. Furthermore, because Flirds
is a Taylor approximation of the per-round Shapley values, the difference
between the two can be interpreted as Taylor approximation error, separated
from differences in game definition or Shapley sampling error. We verify this
relationship by comparing rankings and values against the per-round Shapley
values computed by exhaustive enumeration, and we evaluate training
improvement under a shared intervention policy together with computational
cost. Defining a game that matches the actual FL run, computing it
efficiently, comparing it against the reference obtained by exhaustively
enumerating the same game, and measuring actual training effects to validate
both accuracy and practical effectiveness: this is the structure of the
paper.

Our contributions are as follows.
\begin{enumerate}
    \item \textbf{A Shapley game for the observed federated round.} We
    formalize a validation-loss game whose players are the client updates
    observed in a synchronous FedAvg round, with the actual aggregation
    weights held fixed. Because the grand coalition exactly reproduces the
    actual aggregation, the estimand is explicit, and in-run attribution is
    separated from retraining-based data valuation
    (Sections~\ref{sec:datashapley} and~\ref{sec:game}).
    \item \textbf{A closed-form second-order estimator with interpretable
    error.} We derive the Shapley value of the second-order Taylor
    surrogate of this federated round game. It requires only one
    aggregate-direction server HVP per round and one inner product per
    participant, so contributions can be accumulated online under partial
    participation without per-coalition model evaluation or Shapley sampling;
    a zero-update client receives exactly zero algebraically; efficiency
    holds within the surrogate game; and a player-level error bound follows
    from the Taylor residual
    (Sections~\ref{sec:estimator} and~\ref{sec:theory}).
    \item \textbf{Validation of accuracy and practical effectiveness.} We
    evaluate the ranking and value accuracy of the estimates against the
    per-round Shapley values obtained by exhaustive enumeration. We then
    validate the practical effectiveness of the measured contributions
    through actual training improvement under a shared intervention policy,
    and report measured computational cost for all methods (Section~5).
\end{enumerate}

\section{Related Work}
\label{sec:related}

\textbf{Shapley-based contribution evaluation in federated learning.}
The common bottleneck of this line of work is the exponential number of
coalition evaluations that the Shapley value demands, and
subsequent research has developed in two directions: reducing the number of
coalitions to evaluate, or replacing and complementing the coalition-utility
computation. FedSV \citep{wang2020principled} introduced permutation sampling
and group testing, and GTG-Shapley \citep{liu2022gtg} combined guided Monte
Carlo sampling with truncation within and across rounds on top of submodel
reconstruction from stored updates. ComFedSV \citep{fan2022improving}
completes the utility matrix left partially empty by partial participation
via low-rank matrix completion. SPACE \citep{chen2023space}, in contrast,
introduces single-round knowledge amalgamation and a prototype-based utility,
reducing the cost of multi-round reconstruction and of evaluation that scales
with validation-set size. ShapleyFL and the FedIF and ShapFed families use
Shapley-derived values or gradient similarity for robust aggregation and
client selection
\citep{sun2023shapleyfl,tang2025lightweight,tastan2024redefining}. Ripple
Shapley \citep{zeng2026ripple} decomposes sample-level contribution in a
single federated training run into the direct effect in the current round and
the effects propagated to later rounds, accumulating the latter through a
low-dimensional Jacobian approximation; because its unit of evaluation is the
sample and it targets cross-round influence propagation, its target of
evaluation differs from ours, the client-level contribution in the observed
round game.

\textbf{Data contribution evaluation in centralized training.}
In centralized training, the training party has direct access to individual
samples and their derivative information, and sample-level attribution has
developed on that premise: DataInf, EK-FAC, TRAK, and LoGra estimate
influence and attribution from gradient or curvature information
\citep{kwon2024datainf,grosse2023studying,park2023trak,choe2026your}, and
LESS, MATES, and DsDm use such signals for targeted data selection
\citep{xia2024less,yu2024mates,engstrom2024dsdm}. In-Run Data Shapley (IRDS)
\citep{wang2025data} is the method in this lineage that our paper takes as
its starting point: instead of hypothetical retraining, it fixes a single
actual training run and accumulates the per-step Shapley values of the
validation loss, computed in closed form via Taylor expansion. What an FL server observes, however, is not samples but the
updates submitted by clients, so these methods do not transfer as they are.
What Flirds takes from IRDS is the computational structure in which, once the
per-player model displacements combine additively, the Shapley value of the
second-order surrogate game admits a closed form.

\textbf{Quality evaluation, selection, and markets for federated LLMs.}
FedDQC \citep{du2025feddqc} assesses the quality of a client's internal data
through instruction--response alignment and schedules training accordingly,
and FedHDS \citep{qin2025federated} selects coresets for federated
instruction tuning using redundancy in representation space. iPFL jointly
designs collaboration and payment rules in a personalized-model market. These
studies show that quality control, selection, or incentive mechanisms are
feasible at the scale of federated large language and foundation models, but
they do not compute how much each client's update improved the jointly
trained model in a realized training run.

\section{Problem Formulation and Background on In-Run Attribution}
\label{sec:background}

\subsection{The Observed Federated Learning Run}
\label{sec:fedavg}

In FedAvg \citep{mcmahan2016federated}, clients $k \in [N] := \{1,\dots,N\}$
hold local datasets of size $n_k$, and in round $r$ a participating set
$P_r \subseteq [N]$ performs local training from the current model $w^r$.
Client $k$ sends the resulting model difference $\Delta w_k^r$ to the server,
which aggregates, for $r=0,\dots,R{-}1$,
\begin{equation}
\label{eq:fedavg}
w^{r+1}
= w^r + \sum_{k \in P_r} p_k^r\,\Delta w_k^r,
\qquad
p_k^r = \frac{n_k}{\sum_{j \in P_r} n_j}.
\end{equation}
Here, $p_k^r$ is the weight the server used for the actual participating set
$P_r$.

What the server observes for each client is not raw data but the update
$\Delta w_k^r$, and the model displacement of the actual round is a weighted
sum of these updates. This paper therefore takes the updates submitted by the
clients as the units of evaluation and evaluates each round's contributions
from the observed update vectors and the server's validation loss.

\subsection{Defining Contribution}
\label{sec:datashapley}

Given a player set and a coalition utility, the Shapley value averages each
player's marginal contribution over all participation orders. The value of
player $k$ is
\begin{equation}
\label{eq:shapley}
\phi_k(U)
= \sum_{S \subseteq [N]\setminus\{k\}}
\frac{|S|!\,(N{-}|S|{-}1)!}{N!}
\bigl(U(S\cup\{k\})-U(S)\bigr).
\end{equation}
Within the specified game, this value satisfies efficiency, symmetry, the
null-player property, and linearity \citep{shapley1953value}, but the Shapley
formula itself does not determine which utility is the right one. Client
contribution evaluation must therefore decide, before any computation, which
counterfactual question it is to answer.

Classical Data Shapley uses as utility the final performance of a model
retrained from scratch using only the data in a coalition
\citep{ghorbani2019data}. It asks, ``What model would have been obtained had
only these clients participated in training?'' Per-round contribution in an
actual FL run asks a different question: ``Holding the current model and the
updates actually received fixed, how much did a subset of these updates
contribute to improving this round's model?'' The former is a counterfactual
over complete retraining, whereas the latter is attribution along a realized
trajectory; the two values are answers to different games. We take the latter
as our direct target and analyze its relationship to the former separately
and empirically.

\subsection{The Computational Idea of In-Run Data Shapley}
\label{sec:irds}

In-Run Data Shapley (IRDS) \citep{wang2025data} fixes an actual training run
rather than repeatedly performing hypothetical retraining, and views each
gradient step as a small cooperative game. When step $t$ updates the model as
$w^{t+1}=w^t-\eta_t\sum_{z\in B_t}\nabla\ell_z(w^t)$,
the utility of a batch subset $S\subseteq B_t$ is set to
\begin{equation}
\label{eq:irds-utility}
u_t(S)
= \lval(w^t)
- \lval\left(
w^t-\eta_t\textstyle\sum_{z\in S}\nabla\ell_z(w^t)
\right).
\end{equation}
That is, the decrease in validation loss obtained by applying only the
gradients of some of the samples is taken as the value of that coalition.
Summing each step's Shapley values over the actual training run yields a
per-sample in-run attribution for that trajectory.

The computational key of IRDS is that the per-player model displacements
combine additively. Taylor-expanding the validation loss at the current
model, the first-order term captures the alignment between each player's
displacement and the validation gradient, and the second-order term captures
the curvature effect that arises when a displacement combines with the other
displacements of the same step. Writing
$g_t:=\nabla\lval(w^t)$ and $H_t:=\nabla^2\lval(w^t)$ for the gradient and
Hessian of the validation loss, the Shapley value of the second-order
surrogate is
\begin{multline}
\label{eq:irds-closed}
\phi_z(u_t)
\approx
\eta_t\left\langle g_t,\nabla\ell_z(w^t)\right\rangle\\
-\frac{\eta_t^2}{2}
\left\langle
\nabla\ell_z(w^t),
H_t\textstyle\sum_{z'\in B_t}\nabla\ell_{z'}(w^t)
\right\rangle.
\end{multline}
The important point is that the second-order term of every player can be
computed not from the full Hessian or its inverse, but from a single
Hessian-vector product along the full step direction.

The weighted sum of client updates in FedAvg provides the structure needed to
apply this computational idea at the client level. Simply substituting client
updates for sample gradients, however, does not by itself settle the FL
attribution problem. One must also decide whether the aggregation weights are
held fixed or renormalized within each coalition, how to accumulate values
when the participant set changes from round to round, and which game's
Shapley value the computed quantity targets. The next section first
defines a game that preserves the actual FedAvg displacement and then
explains how Flirds reduces the cost of computing that game.

\section{Flirds: The Federated In-Run Game and Its Closed-Form Estimation}
\label{sec:flirds}

\subsection{The Federated Game and the Two Ground Truths}
\label{sec:game}

The classical ground truth of data valuation, the \emph{exact retrain
Shapley}, retrains from scratch for $R$ rounds with only the client subset $S$
participating and takes the validation-loss decrease of the resulting final
model $w_S^R$ (with $w_\emptyset^R = w^0$),
\begin{equation}\label{eq:retrain}
U^{\mathrm{re}}(S) = \lval(w^0) - \lval\big(w_S^R\big),
\end{equation}
whose Shapley value $\phi_k^{\mathrm{re}} = \phi_k(U^{\mathrm{re}})$ is
computed without approximation by retraining all $2^N$ subsets (a
method-neutral ground truth; the question it asks is the counterfactual ``had
only $S$ participated''). The round structure of FL, meanwhile, corresponds
exactly to the step structure of IRDS: in the step game of
Eq.~\eqref{eq:irds-utility}, substituting the round's weighted per-client
update $p_k^r \Delta w_k^r$ for the sample gradient term
$-\eta_t \nabla \ell_z$ transplants the game to the client level, and the
ingredients are nothing but what the server receives anyway under the standard
protocol. We freeze the realized trajectory
$\{w^r, \{\Delta w_k^r\}_{k \in P_r}\}_{r=0}^{R-1}$ (the \emph{frozen
trajectory}; every method below consumes this same training log) and define
the per-round coalition utility, for $S \subseteq P_r$, as
\begin{equation}\label{eq:round-utility}
u_r(S) = \lval(w^r)
- \lval\Big(w^r + \textstyle\sum_{k \in S} p_k^r\, \Delta w_k^r\Big).
\end{equation}
The weights are the \emph{fixed weights} $p_k^r$ that the run itself used; we
do not renormalize within $S$. Renormalizing within $S$, as some baselines do
(GTG-Shapley, FedSV), is not an implementation detail but a \emph{different
cooperative game}: a zero-update free-rider comes to receive a nonzero value;
there are counterexamples in which rankings flip under unequal data sizes,
partial participation, and second-order curvature (Appendix~A); and the
sum-form structure that folds a round into a single HVP breaks, so the
closed-form computation no longer holds. The \emph{exact in-run Shapley}
$\phi^{\mathrm{in}}$ is the sum of each round game's Shapley value computed
without approximation by exhaustively enumerating all $2^{|P_r|}$ subsets per
round: $\phi_k^{\mathrm{in}} = \sum_{r:\, k \in P_r} \phi_k(u_r)$. The
efficiency of IRDS carries over, combined with telescoping:
$\sum_k \phi_k^{\mathrm{in}} = \sum_{r=0}^{R-1} u_r(P_r)
= \lval(w^0) - \lval(w^R)$.

\textbf{Terminology.} Throughout, ground truth (oracle) refers to exactly two
objects: the \emph{exact retrain Shapley} above (exhaustive per-subset
retraining) and the \emph{exact in-run Shapley} just defined (exhaustive
per-round enumeration). We always use the full names and never an unqualified
``exact.'' For both, larger is more beneficial (positive when the validation
loss is reduced).

\subsection{The Closed-Form Estimator}
\label{sec:estimator}

The exact in-run Shapley involves no approximation, but it demands
$2^{|P_r|}$ validation evaluations per round. Flirds starts from the fact that
the Shapley value of $\hat u_r$, the approximate game obtained by
Taylor-expanding the round utility to second order around $w^r$, comes out in
closed form (Proposition~\ref{prop:closed}). Writing
$g_r := \nabla \lval(w^r)$, $H_r := \nabla^2 \lval(w^r)$ (the true Hessian,
not a Gauss--Newton approximation), $\delta_k^r := \Delta w_k^r$, the
aggregate displacement of a coalition
$\Delta_S^r := \sum_{k \in S} p_k^r\, \delta_k^r$, and
$\Delta W_r := \Delta_{P_r}$:
\begin{equation}\label{eq:closed}
\begin{split}
\hat\phi_k^{(r)}
&= -\,p_k^r\, \big\langle g_r,\, \delta_k^r \big\rangle
- \tfrac{1}{2}\, p_k^r\, \big\langle \delta_k^r,\; H_r\, \Delta W_r
\big\rangle, \\
\hat\phi_k &= \sum_{r\,:\,k\in P_r} \hat\phi_k^{(r)}.
\end{split}
\end{equation}
The first term measures how well $k$'s update aligns with the descent
direction of the validation loss; the second, the \emph{client-interaction
term}, measures how $k$'s update composes, through curvature, with the other
participants' aggregate update. (The sign of the first term differs from
Eq.~\eqref{eq:irds-closed} because there the displacement
$-\eta_t \nabla \ell_z$ carries an explicit minus sign, whereas here
$\delta_k^r$ is itself the displacement.) All a round requires is \emph{one}
Hessian-vector product (HVP) $H_r \Delta W_r$, computed with forward-mode
automatic differentiation, plus one inner product per participant. No Hessian
is ever materialized or inverted, and the cost is independent of the number of
subsets $2^{|P_r|}$. The variant with the second-order term removed is called
\emph{Flirds (first-order)} and serves as an ablation of the interaction term.

\subsection{Theory: The Estimator and the Exact In-Run Shapley Compute the
Same Game}
\label{sec:theory}

The point of this section is one sentence: Flirds computes not an arbitrary
score but the Shapley value of precisely the game defined by
Eq.~\eqref{eq:round-utility}, and its sole difference from the exact in-run
Shapley is the Taylor truncation. Formal statements, proofs, assumptions,
counterexamples, and numerical verification are in Appendix~A.

\begin{proposition}[Per-round decomposition]\label{prop:decomp}
On the frozen trajectory, the Shapley value of the total game $\sum_r u_r$
(each $u_r$ extended to a game on $[N]$ by $u_r(S) := u_r(S \cap P_r)$ before
summing) decomposes exactly into the sum of the Shapley values of the
per-round $|P_r|$-player games. (That non-participating rounds contribute
zero is a theorem, not a definition.)
\end{proposition}

\begin{proposition}[Closed form; free-rider zero; efficiency]
\label{prop:closed}
For the second-order Taylor approximation game $\hat u_r$ of the round
utility, the exact Shapley value equals Eq.~\eqref{eq:closed} exactly:
$\hat\phi_k^{(r)} = \phi_k(\hat u_r)$ (proof in Appendix~A). Two corollaries:
a zero-update free-rider with $\delta_k^r = 0$ in every participating round
receives \emph{exactly} $\hat\phi_k = 0$ algebraically, and the values sum to
the total utility of the approximate game (efficiency).
\end{proposition}

\begin{proposition}[Taylor residual]\label{prop:residual}
The price of the approximation is truncation error. The per-round utility
error is bounded by $\frac{M_3^r}{6}\|\Delta_S^r\|^3$ at second order and by
$\frac{M_2^r}{2}\|\Delta_S^r\|^2$ at first order, and accumulates linearly in
the number of rounds (the per-round local curvature constants
$M_2^r, M_3^r$ are defined formally in Appendix~A). These constants are not
estimable quantities, so what the bound delivers is the \emph{order} of the
residual (its scaling law).
\end{proposition}

\begin{remark}[Scope]
The exact in-run Shapley is a function of the realized trajectory, and
therefore inherits IRDS's unresolved axiomatization of trajectory-specific
values (its relation to the exact retrain Shapley is an empirical question
rather than a theoretical one, and Section~5.3 reports it as such). The game
is defined at \emph{client} granularity, the unit the FL protocol exposes to
the server; the formal per-sample-to-per-client bridge holds under mean losses
and a single local step (Appendix~A) and does not transfer as-is to token-mean
LLM losses. In the LLM track, $w$, $\delta_k^r$, $g_r$, and $H_r$ are all
quantities in the LoRA-factor coordinates that the server actually exchanges
and averages (the CNN track trains full parameters, so no distinction arises;
the formal treatment of coordinate dependence is in Appendix~A.10).
\end{remark}

% \section{Experiments}
%   ← §5 실험 (5.1 세팅·측정 프로토콜 / 5.2 추정 정확도 / 5.3 참값 대비·게임-무관 척도 / 5.4 신호 조건 / 5.5 비용 / 5.6 기여도의 가치: 부호-게이팅·selection·탐지)

% \section{Discussion and Limitations}
%   ← §6 논의와 한계

% \section{Conclusion}
%   ← §7 결론

% (부록 A--F: 학회 규정에 따라 본문 뒤 technical appendix 또는 별도 제출 — paper-ko.md 부록 참조)

% References (bib 파일: 같은 폴더 references.bib — Overleaf는 프로젝트 루트 밖(../)을
% 읽지 못하므로 main.tex 옆에 둔다; \bibliographystyle은 aaai2027.sty가 내부 설정)
\bibliography{references}

% Check whether the conference requires a reproducibility checklist to be included in the paper.
% If so, you can uncomment the following line and ajust the path to include it.
% \input{ReproducibilityChecklist.tex}

\end{document}
